\section{问题一：静水风载下的系统响应}
\subsection{本问分析}
问题一给定II型锚链22.05 m、重物球1200 kg、水深18 m和静水条件，要求在12 m/s与24 m/s风速下计算浮标吃水、各刚性构件倾角、锚链形状和游动区域。其耦合关系为：吃水决定露出水面的迎风面积，风力决定水平张力；下部构件和悬起锚链的竖直载荷决定浮标吃水；构件姿态与锚链高度共同满足18 m水深闭合。

\subsection{浮标与刚性构件平衡}
设浮标吃水为 \(h\)，其浮力和风荷载分别为

\[
B_f=\rho g\pi h,
\qquad
H=1.25(2-h)v^2. \tag{1}
\]
钢管、钢桶按封闭圆柱的外形体积计算浮力，每节钢管和钢桶的水中有效重量分别为

\[
W_p=78.3566\ \mathrm N,
\qquad W_b=270.2363\ \mathrm N.
\]
题目没有提供重物球体积，故将其作为钢桶下端的独立集中载荷并忽略浮力，取

\[
W_s=1200g=11772\ \mathrm N.
\]
对于任一刚性构件，设上下端竖直张力为 \(V_u,V_l\)，则

\[
V_u=V_l+W,
\qquad
\tan\theta=\frac{2H}{V_u+V_l}. \tag{2}
\]
锚链顶端竖直张力为 \(V_t\) 时，钢桶下端载荷为

\[
V_{b,l}=V_t+W_s.
\]
由钢桶向上逐节递推四节钢管，得到下部系统对浮标的竖直拉力

\[
V_0=V_t+W_s+W_b+4W_p. \tag{3}
\]
浮标竖直平衡给出

\[
h=\frac{1000g+V_0}{\rho g\pi}. \tag{4}
\]
式（1）和式（4）形成吃水与水平风载之间的反馈。

\subsection{锚链形状与状态判断}
II型锚链单位长度质量为7 kg/m，忽略锚链浮力后

\[
q=7g=68.67\ \mathrm{N/m}.
\]
若存在铺底段，设悬起链长为 \(l\)，离床点处切线水平，则

\[
V_t=ql,
\]
\[
x_c=\frac{H}{q}\operatorname{arsinh}\frac{ql}{H},
\]
\[
y_c=\frac{H}{q}
\left[
\sqrt{1+\left(\frac{ql}{H}\right)^2}-1
\right]. \tag{5}
\]
水深闭合条件为

\[
h+y_c+\cos\theta_b+
\sum_{i=1}^{4}\cos\theta_i=18. \tag{6}
\]
由式（6）求出 \(l\)。若 \(l\le22.05\)，铺底假设成立；若超过总链长，则转入全悬链模型。本问两种风速求得的悬起链长均小于总长，因此锚点处切线水平，锚点角为零。

浮标游动半径由铺底链段、悬起链段及五个刚性构件的水平投影组成：

\[
R=(22.05-l)+x_c+\sin\theta_b+
\sum_{i=1}^{4}\sin\theta_i. \tag{7}
\]
游动圆面积取 \(A=\pi R^2\)。

\subsection{求解结果}
\begin{table}[H]\centering\small
\caption{问题一两种风速下的系统平衡状态}
\label{tab:q1-1}
\begin{tabular}{ccc}
\toprule
指标 & 12 m/s & 24 m/s \\
\midrule
吃水深度/m & 0.7348 & 0.7489 \\
第1节钢管倾角/(°) & 0.9764 & 3.7324 \\
第2节钢管倾角/(°) & 0.9822 & 3.7536 \\
第3节钢管倾角/(°) & 0.9880 & 3.7751 \\
第4节钢管倾角/(°) & 0.9939 & 3.7968 \\
钢桶倾角/(°) & 1.0073 & 3.8461 \\
悬起链长/m & 15.2254 & 21.7267 \\
铺底链长/m & 6.8246 & 0.3233 \\
游动半径/m & 14.3029 & 17.4232 \\
游动圆面积/m² & 642.68 & 953.68 \\
\bottomrule
\end{tabular}
\end{table}
风速加倍后，风荷载近似增至4倍，构件倾角和游动半径明显增大，铺底长度由6.8246 m减小至0.3233 m；浮标吃水仅增加0.0141 m，说明风速主要改变水平载荷。24 m/s工况已接近锚链完全离床临界状态，但铺底长度仍为正，采用部分铺底模型自洽。

两种工况的链形可由式（5）直接参数化。12 m/s时$H/q=3.3164$ m，铺底段终点位于$x=6.8246$ m，悬链段延伸至$x=14.2165$ m并升高12.2660 m；24 m/s时$H/q=13.1176$ m，铺底段终点移至$x=0.3233$ m，悬链段延伸至$x=17.0935$ m并升高12.2620 m。风速增加使特征长度$H/q$明显增大，链形趋于平缓并向水平方向展开，而顶端高度因18 m水深闭合保持接近。该参数变化与图\ref{fig:q1-chain}所示曲线一致，也解释了游动半径增加而吃水变化较小的原因。

\subsection{连续模型验证}
II型锚链单环长度为0.105 m，总长对应210个链环。将锚链离散为短柔性单元，以单元两端竖直张力平均值确定其方向，并与浮标、刚性构件和水深方程联立。连续与离散模型在12 m/s、24 m/s下的游动半径相对误差分别为0.000818\%和0.000145\%，钢桶与钢管倾角相对误差均小于0.0001\%。最大相对误差出现在24 m/s的铺底长度，为0.010738\%，但绝对差仅 \(3.47\times10^{-5}\) m。由此说明连续悬链线足以用于后续配重搜索和多工况设计。

校核数据共比较两种风速下的吃水、水平张力、悬起链长、铺底链长、锚链跨度与高度、五个构件倾角及游动半径等24项指标。12 m/s工况的最大相对误差为0.002043\%，对应铺底长度；24 m/s工况的最大相对误差为0.010738\%，同样对应铺底长度。铺底长度是总链长与悬起链长之差，在24 m/s时仅为0.3233 m，较小的绝对误差除以接近零的铺底长度会放大相对误差，因此不能据此判断连续模型失效。更直接影响结论的游动半径、吃水和钢桶倾角均保持更高一致性，且两种模型对“仍有铺底链段”的状态判定相同。

连续模型与离散模型的误差方向也具有可解释性。链环模型用有限个线段逼近光滑悬链线，节点处方向离散会使水平跨度和离床位置产生微小偏差；随着链环数量增加，这类几何误差应收敛。由于本题已由实际210个链环完成离散，比较并非任意网格试算，而是对应题设结构的独立复算。校核的目的在于确认连续悬链线不会改变主要响应和状态边界，不把远低于测量尺度的数值差异包装成工程精度。

\subsection{求解流程、结果解释与适用范围}
给定悬起链长后，程序由锚链顶端向上递推钢桶、四节钢管及浮标的竖直张力和倾角，并同步更新吃水与风荷载；外层用二分法使总竖直高度等于18 m。若悬起长度超出总链长，才切换到完全悬起模型。本问两种风速均落在部分铺底分支，并通过受力、力矩和水深闭合检查。

风速由12 m/s增至24 m/s后，悬起链长增加6.5013 m，而吃水仅增加0.0141 m，说明游动范围和铺底长度比吃水更敏感。24 m/s时仅余0.3233 m锚链铺底，提示更高风速下必须允许状态切换。210链环离散复算与连续模型的关键响应相对误差均小于0.011\%，支持后续批量计算采用连续悬链线。

模型只描述稳定风作用下的静态平衡，不能刻画结构惯性、阵风峰值和链环接触滞回；理想柔索近似也忽略轴向伸长与弯曲刚度。离散校核仅证明连续悬链线在当前稳态工况下的一致性，不代表动态响应和疲劳安全性。

\subsection{本问结论}
在静水和给定结构条件下，12 m/s与24 m/s风速对应的钢桶倾角分别为 \(1.0073^\circ\) 与 \(3.8461^\circ\)，浮标吃水分别为0.7348 m与0.7489 m，游动半径分别为14.3029 m与17.4232 m。两种工况锚链均有铺底段，锚点夹角为零；24 m/s工况更接近整链悬起，是进入问题二极端风速分析的基础。
\aimark{3}
